\section{Introduction}

Let $h>0$ and let $(v_x)_{x\in\Z^2}$ be independent random variables,
each uniform on $[-h,h]$. We consider the bounded self-adjoint operator
\begin{equation}\label{setup:model}
 (H_v u)(x)=\sum_{\abs{y-x}_1=1}u(y)+v_xu(x)
 \qquad\text{on }\ell^2(\Z^2).
\end{equation}
Thus $v_x=h\omega_x$ in the usual notation with $\omega_x$ uniform on
$[-1,1]$. Pure-point spectral type means that the closed span of the
eigenvectors is the entire Hilbert space; it does not mean that every
point of the spectrum is an eigenvalue.

The planar spectral localization problem asks whether this pure-point
conclusion holds throughout the spectrum for every nonzero width of the
uniform distribution \cite[Example~2.1 and Problem~2]{Simon2000}.
We prove that it does.

\begin{theorem}\label{main:theorem}
For each fixed $h>0$, almost surely $H_v$ has pure-point spectral type,
admits a complete orthonormal eigenbasis, and has spectrum
\[
 \sigma(H_v)=[-4-h,4+h].
\]
Its eigenvalues are dense in this interval. The probability-one event
in this statement is allowed to depend on $h$.
\end{theorem}

The conclusion includes arbitrarily weak nonzero disorder and the entire
spectrum. It resolves the pure-point assertion of the two-dimensional
Anderson localization conjecture for a uniform single-site distribution,
as formulated in \cite[Problem~2]{Simon2000}. An arbitrary nondegenerate
uniform interval reduces to the centered interval above by subtracting
its midpoint from the operator. We do not assert one probability-one
event for all disorder strengths, a bound uniform in energy, or a
dynamical-localization theorem.

\subsection*{Context and scope}
Anderson introduced random site energies as a mechanism for suppressing
quantum transport, motivated by spin diffusion and impurity conduction
